\documentclass[11pt]{article}

\usepackage[margin=1in]{geometry}
\usepackage[T1]{fontenc}
\usepackage{lmodern}
\usepackage{amsmath,amssymb,mathtools}
\usepackage{booktabs,tabularx,longtable,array}
\usepackage[numbers,sort&compress]{natbib}
\usepackage{microtype}
\usepackage{xcolor}
\usepackage{hyperref}
\usepackage{lineno}
\usepackage{setspace}

\definecolor{icpnavy}{HTML}{183B56}
\definecolor{icpteal}{HTML}{2B6F77}

\hypersetup{
  colorlinks=true,
  linkcolor=icpnavy,
  citecolor=icpteal,
  urlcolor=icpteal,
  pdftitle={The Constrained Future: A Testable Information-Constraint Principle for Persistent Organization},
  pdfauthor={Douglas Clark},
  pdfsubject={BioSystems hypothesis paper, revision 1},
  pdfkeywords={information constraint, adjacent possible, constraint closure, epigenetic memory, carrier succession, persistent organization, relative entropy}
}

\setlength{\tabcolsep}{5pt}
\renewcommand{\arraystretch}{1.14}
\emergencystretch=2em
\onehalfspacing

\newcommand{\ICP}{\textsc{ICP}}
\newcommand{\IOTA}{\textsc{IOTA}}
\newcommand{\Efamily}{\mathcal{E}_{\tau}}
\newcommand{\Rset}{\mathcal{R}}
\newcommand{\status}[1]{\textsf{#1}}

\title{\textbf{The Constrained Future}\\
\large A Testable Information-Constraint Principle for Persistent Organization}
\author{Douglas Clark\\
\small Independent Researcher\\
\small Corresponding author: \href{mailto:clarkdd@gmail.com}{clarkdd@gmail.com}}
\date{}

\begin{document}
\maketitle

\noindent\textbf{Article type:} Hypothesis Paper

\begin{abstract}
Living organization persists while its material components are replaced. We
formulate the Information-Constraint Principle (ICP) as an operational
hypothesis about such persistence: physically effective information is an
interaction-formed, embodied constraint on an explicitly indexed family of
reachable or probable future responses, and it persists only when its causal
role is re-instantiated in successor carriers. A critical comparison with
viability theory, the adjacent possible, empowerment, counterfactual semantic
information, computational mechanics, constructor theory, and feedback
thermodynamics shows that most future-option language is established prior
art. ICP's remaining content is a persistent-carrier test requiring formation,
sufficiency, carrier intervention, succession across replacement, and
representation robustness. We formalize the adjacent possible as the
novelty-bearing projection of an indexed reachable set and show that a signed
support deformation equals forward Kullback--Leibler divergence under pruning
and negative reverse divergence under enabling. Existing models of the adjacent
possible postulate the step that opens new possibilities; ICP makes that step
an intervention on an embodied carrier and predicts, unlike triggering models,
that opened possibilities close when their carrier is removed, as permissive
substitutions in glucocorticoid-receptor evolution illustrate. Applied to yeast \textit{INO1}
transcriptional memory, the test shows that a chromatin record survives several
rounds of replication-coupled carrier replacement while depending on a
maintenance system, but that succession of the functional role has not been
tested; we specify the experiment that would test it and the
nucleosome-modification model it must be compared against. ICP is a
falsifiable organizing hypothesis, not an established law: its decisive test is
whether carrier interventions predict held-out changes in future responses
beyond complete native dynamics.
\end{abstract}

\noindent\textbf{Keywords:} information constraint; adjacent possible;
constraint closure; epigenetic memory; carrier succession; persistent
organization

\linenumbers

\section{The persistence problem}

Living organization persists while its material components are replaced.
Enzymes are degraded and resynthesized, nucleosomes are diluted at every
replication fork, and cells in a tissue turn over while the tissue continues to
perform its role. The motivating question is how the consequence of an
interaction can persist after the object or arrangement that first carried it
has changed or disappeared, without slipping between incompatible meanings of
``information.'' None of these examples requires information to float free of
physical realization. They distinguish persistence beyond an \emph{original}
carrier from existence without \emph{any} carrier.

Embodied-information accounts in theoretical biology already treat patterns as
constraints on physical processes and emphasize their historical accumulation
\citep{Farnsworth2022Information}. Constraint-closure theory models conserved
structures that act on processes while depending on one another for their
maintenance \citep{MontevilMossio2015Closure}, and the theory of the adjacent
possible treats organization as enabling futures that were previously
inaccessible \citep{Kauffman2000Investigations,Kauffman2014Prolegomenon}. A new
proposal must therefore earn its place through operational criteria, not
through the language of embodiment or constraint alone.

We state the proposal as follows:
\begin{quote}
\textbf{Information-Constraint Principle (ICP).} Physically effective
information is interaction-formed, embodied constraint on an explicitly indexed
family of reachable or probable futures; it persists when that constraint role
is lawfully re-instantiated beyond its original carrier.
\end{quote}
ICP arose within a broader information-first modelling program whose
artificial-life testbed, IOTA, is reported separately
\citep{Clark2026IotaConstraint}. The present paper does not depend on that
program's physical ontology. It asks a narrower question: does the combination
of historical formation, physical embodiment, sufficiency for a declared
family of future responses, causal carrier manipulation, succession across
replacement of the original carrier, and robustness to representation yield a
testable claim that existing theories do not already make?

The phrase ``constrains the future'' is easy to satisfy. Given a complete
physical state and law, virtually every state difference can alter a future
trajectory. Control theory studies reachable sets; computational mechanics
conditions futures on pasts; and semantic-information work erases correlations
and measures effects on later viability. If ICP merely renames these objects it
may be useful language, but it is not a new principle. The paper is
accordingly adversarial: it formalizes ICP (Section~\ref{sec:formalism}), tests
it against its closest native formalisms (Section~\ref{sec:redundancy}),
applies it to a well-characterized biological memory
(Sections~\ref{sec:biology}--\ref{sec:ino1}), and grades what survives
(Sections~\ref{sec:audit}--\ref{sec:path}).

\section{Formalism}
\label{sec:formalism}

This article is a critical integrative synthesis and operational proposal, not
a systematic review. A project registry records query strata, primary-source
checks, source notes, competing explanations, falsification criteria, and a
claim--evidence matrix; searches were expanded adversarially whenever a field
offered a likely redundancy account. Claims are labelled as established field
results, project evidence, synthesis inferences, or ICP hypotheses, and project
papers are never counted as independent replication.

\subsection{Indexed future-law families}

For a candidate retained state or record $m$, intervention family $\Pi$,
observable algebra $\mathcal A$, boundary and resolution declaration $B$,
resource declaration $C$, and horizon $\tau$, define
\begin{equation}
 \Efamily(m;\Pi,\mathcal A,B,C)
 =\left\{P^{\pi}_{\mathcal A}
   \bigl(F_{t:t+\tau}\mid m,B,C\bigr):\pi\in\Pi\right\}.
 \label{eq:future-family}
\end{equation}
The corresponding reachable support for policy $\pi$ is
\begin{equation}
 \Rset^{\pi}_{\tau}(m)=
 \operatorname{supp}P^{\pi}_{\mathcal A}
 \bigl(F_{t:t+\tau}\mid m,B,C\bigr),
 \label{eq:reachable-support}
\end{equation}
with topology and tolerance specified for continuous spaces.

Equation~\eqref{eq:future-family} is not itself an ICP law; it is shared
comparison language. A \emph{constraint} requires a matched contrast between
families, for example a history-formed carrier $m$ and a reference $m_0$
prepared with the same declared ordinary state variables and resources:
\begin{itemize}
  \item a \textbf{hard constraint} changes reachable support on a common path
  or outcome space;
  \item a \textbf{soft constraint} changes weights while support may remain
  fixed; and
  \item an \textbf{operational constraint} changes a reachability or
  convertibility preorder under the same allowed operations.
\end{itemize}

\subsection{A relative-entropy constraint diagnostic}

For a declared comparison $m$ versus $m_0$ at $x$, one distribution-sensitive
member of the future-family comparison is
\begin{equation}
 \mathfrak C_{x,\tau}^{\pi}
 (m\Vert m_0;\mathcal A,B,C)
 :=D_{\mathrm{KL}}\!\left(
 P^{\pi}_{\mathcal A}(F_{t:t+\tau}\mid m,B,C)
 \middle\Vert
 P^{\pi}_{\mathcal A}(F_{t:t+\tau}\mid m_0,B,C)
 \right),
 \label{eq:icp-constraint-divergence}
\end{equation}
provided both laws use the same outcome space, intervention, horizon, measure,
and coarse-graining, with the first law absolutely continuous with respect to
the second \citep{KullbackLeibler1951Information}. This is an indexed
diagnostic, not the definition of ICP: the full family retains outcome
identity, topology, interventions, and operational order that a single scalar
discards. For a soft constraint it is the natural quantity; for a hard
constraint it has an exact support specialization. If the post-intervention
law is uniform on $\Rset_x^{\mathrm{post}}\subseteq\Rset_x^{\mathrm{prior}}$
and the reference law is uniform on $\Rset_x^{\mathrm{prior}}$, then
\begin{equation}
 D_{\mathrm{KL}}(U_x^{\mathrm{post}}\Vert U_x^{\mathrm{prior}})
 =\log\frac{|\Rset_x^{\mathrm{prior}}|}{|\Rset_x^{\mathrm{post}}|}.
 \label{eq:iota-kl-identity}
\end{equation}
The direction matters: the reversed divergence is infinite under strict
support pruning.

\subsection{The adjacent frontier and signed possibility-space flux}

The adjacent possible becomes operational only after its outcome types and
notion of ``one step'' are declared. Let $r_\ell:\Omega\rightarrow\mathcal
K_\ell$ be a frozen scale-$\ell$ map from micro-outcomes to registered types or
roles, and let $\mathcal V_t^\ell$ be the types actualized by time $t$. Define
\begin{equation}
 \mathcal J_{t,\tau}^{\pi,\ell}(m_t)
 :=r_\ell\!\left(\Rset_\tau^\pi(m_t)\right)\setminus\mathcal V_t^\ell .
 \label{eq:adjacent-frontier}
\end{equation}
When $\tau$ is one registered transition, this is the novelty-bearing sector of
Kauffman's adjacent possible: types not yet actual but reachable in one step
from the actual \citep[p.~142]{Kauffman2000Investigations}. It depends on
$r_\ell,\pi,\tau,B,$ and $C$; changing a classifier, resource budget, or policy
can change the set.

For matched pre/post supports $R_0$ and $R_1$, retain the identity of the
change rather than only its cardinality:
\begin{equation}
 K=R_0\cap R_1,\qquad E=R_1\setminus R_0,\qquad L=R_0\setminus R_1,
 \label{eq:possibility-flux-ledger}
\end{equation}
the retained, enabled, and lost possibilities. A signed log-volume deformation
is $\chi(R_1:R_0):=\log[\mu(R_0)/\mu(R_1)]$. For nonempty finite uniform nested
supports it gives the exact dual pair
\begin{equation}
 \chi(R_1:R_0)=
 \begin{cases}
 D_{\mathrm{KL}}(U_{R_1}\Vert U_{R_0}),&R_1\subseteq R_0
 \quad\text{(pruning)},\\[3pt]
 -D_{\mathrm{KL}}(U_{R_0}\Vert U_{R_1}),&R_0\subseteq R_1
 \quad\text{(enabling)}.
 \end{cases}
 \label{eq:signed-kl-duality}
\end{equation}
Thus $\chi>0$ denotes net contraction and $\chi<0$ net expansion. For
non-nested supports $\chi$ is only a net-volume contrast and can vanish under
complete identity turnover, so the ledger in
Eq.~\eqref{eq:possibility-flux-ledger} must be reported alongside it.

This duality makes one biologically familiar pattern precise. A retained
carrier may restrict immediate microstate options while enabling more
long-horizon role types,
\begin{equation}
 \chi_{\mathrm{micro},1}>0,\qquad
 \chi_{\mathrm{role},h}<0\quad(h\ge h_\star),
 \label{eq:constraint-enablement-crossover}
\end{equation}
under the same randomized carrier intervention. Enabling-constraint and
scaffolding accounts already anticipate this qualitative pattern
\citep{Kauffman2014Prolegomenon,Griesemer2024Scaffolded}; it becomes an ICP
result only if $h_\star$ is predicted from independently measured carrier
lifetimes and confirmed on held-out systems.

\subsection{The persistent-carrier criterion}
\label{sec:criterion}

A proposed informational carrier $M$ passes the ICP criterion only if four
gates are met:
\begin{description}
  \item[Formation.] Controlled histories create the candidate distinction.
  \item[Sufficiency.] The distinction is sufficient for the declared future-law
  family to a frozen error tolerance.
  \item[Carrier intervention.] Ablating, replacing, transplanting, or scrambling
  the proposed carrier changes a preregistered component of the future family,
  while matched shams and resource controls do not.
  \item[Robustness.] The result survives adequate state augmentation and
  reasonable representation, partition, and coarse-graining changes, or its
  representation dependence is stated as a boundary.
\end{description}
Persistence across carrier succession adds one further demand. Let
$c_0,c_1,\ldots$ be physical realizations at successive times. They instantiate
the same role only relative to a declared equivalence relation
$\sim_{\Pi,\mathcal A,B,C,\tau}$ under which replacement preserves the relevant
future-response family. Material disjointness alone does not establish
sameness of role; sameness of an analyst's label does not establish physical
transport.

\subsection{Novelty and failure tests}

The strongest null is complete-state ordinary dynamics: once every causally
relevant variable, law, input, boundary, and control is included, there is no
additional historical variable. ICP earns distinct content only by delivering
at least one of the following beyond the best native model: a minimal
sufficient physically embodied equivalence class; a representation-stable
carrier-succession invariant; a new impossibility or conversion result; a
demonstrably superior compression of intervention--response families; or a
novel preregistered prediction. Disappearance under adequate state
augmentation, failure of matched carrier interventions, or dependence on one
convenient representation would lower or defeat the relevant claim.

\section{Closest native formalisms and the redundancy test}
\label{sec:redundancy}

The literature that matters most for ICP is the literature capable of making
it redundant.

\subsection{Reachability and viability}

Hamilton--Jacobi reachability characterizes finite-horizon backward-reachable
sets of nonlinear differential games as sublevel sets of viscosity solutions
\citep{MitchellEtAl2005Reachable}, and viability theory identifies states from
which some admissible evolution remains inside a declared constraint set
\citep{SaintPierre1994Viability}. These are mature theories of accessible and
survivable futures. A reachability calculation is, however, given the state,
law, controls, and horizon; it need not ask how a constraint role was formed,
where it is physically realized, or whether it survives replacement of its
first carrier. That upstream question is the only place ICP can add content.

\subsection{The adjacent possible}

Kauffman's chemical adjacent possible---species absent from the actual but
synthesizable from it in one step \citep[p.~142]{Kauffman2000Investigations}---is
exactly the type-projected one-step frontier of
Eq.~\eqref{eq:adjacent-frontier}. His BioSystems formulation adds that enabling
constraints create unprestatable possibilities in evolution
\citep{Kauffman2014Prolegomenon}, and a recent extension contrasts law-governed
possibility spaces with a ``domain of no law'' \citep{PatraKauffman2025NoLaw}.
Fari\~na and Vallejo define a biological cosmological constant proportional to
the squared fractional expansion rate of Fisher-information volume and read the
origin of life as a geometric phase transition \citep{FarinaVallejo2026LambdaB};
this is explicit prior art for geometric expansion into the adjacent possible,
although it does not identify a replaceable carrier or distinguish signed
pruning from enabling.

Formal models make the idea concrete in non-equivalent ways: urns with
triggering yield correlated novelty with Heaps- and Zipf-like scaling
\citep{TriaEtAl2014Novelties}; reinforced walks place discovery on concept
networks \citep{IacopiniEtAl2018Innovation}; scaffolding changes which
possibilities are developmentally available \citep{Griesemer2024Scaffolded};
and autocatalytic closure can causally change complexity dynamics in a Kauffman
chemical model \citep{Faggian2025Closure}. Combinatorial TAP models produce long
plateaus followed by explosive growth \citep{SteelEtAl2020Combinatorial}, while
empirical product data show resource-constrained rather than unrestricted
combinatorial discovery \citep{Taalbi2026AdjacentPossible}.

ICP therefore cannot claim history-dependent accessibility, triggered support
growth, scaffolded opportunity, a plateau-to-explosion transition, or the first
geometric measure of adjacent-possible expansion. The remaining addition is
causal: predeclare the type map, frontier, carrier, horizon, and resources;
alter or replace the carrier; and predict the enabled, lost, and retained
frontier on held-out interventions. The claim that biological phase space is
intrinsically unprestatable \citep{LongoEtAl2012Enablement} cannot
simultaneously serve as an enumerated experimental variable without an
operational surrogate.

\subsection{Empowerment and semantic information}

Empowerment is the channel capacity from action sequences to later sensed
states \citep{KlyubinEtAl2005Empowerment}, and causal entropic forcing builds a
path-entropy objective into an effective force
\citep{WissnerGrossFreer2013CausalEntropy}. These defeat any claim that ICP
discovered the valuing of future options. ICP does not assert that systems
maximize optionality: a constraint may reduce, redirect, or reorganize
futures, and scalar option values collapse outcome identity.

Kolchinsky and Wolpert provide the closest formal rival. They define semantic
information as system--environment correlation whose controlled scrambling
lowers a declared future viability \citep{KolchinskyWolpert2018Semantic}; the
intervention, boundary, horizon, and viability function are constitutive.
Causal-information measures likewise require post-intervention distributions
\citep{AyPolani2008InformationFlows}, and predictive association alone is not a
carrier test \citep{JamesEtAl2016InformationFlows}. ICP is broader only in that
its target may be any registered future-response role rather than viability,
and in its emphasis on formation and succession. On a self-maintenance task an
ICP analysis should specialize to semantic information, not rename it.

\subsection{Computational mechanics and constructor theory}

Causal states are equivalence classes of pasts with the same conditional
distribution over futures and give a minimal optimal predictive representation
\citep{ShaliziCrutchfield2001ComputationalMechanics}; epsilon-transducers extend
this to input-dependent futures \citep{BarnettCrutchfield2015Transducer}. This
is almost exactly the $H\rightarrow M\rightarrow P(F\mid M)$ part of ICP. The
difference is physical: an inferred causal state need not be localized,
manipulable, or interventionally defined, whereas the carrier criterion asks
whether an embodied realization is necessary and sufficient and whether its
replacement preserves the role. Constructor theory supplies a law-level
language of possible and impossible tasks, including information as copyable
and distinguishable variables \citep{Deutsch2013Constructor,
DeutschMarletto2015Information}, but not a history-conditioned process.

\subsection{Feedback thermodynamics}

Feedback thermodynamics shows that a retained measurement changes later work
bounds \citep{SagawaUeda2010Feedback}, demonstrated experimentally with a
Brownian particle \citep{ToyabeEtAl2010InformationEnergy} and made fully
autonomous in a bit-tape demon \citep{MandalJarzynski2012Demon}. It supports
causal utility and closed accounting of memory, but energy, baths, and
protocols enter as supplied structure.

\subsection{Result of the audit}

The generic proposition that present organization constrains the future is
preempted. The viable target is narrower:
\begin{quote}
An interaction forms a physically embodied constraint role; the original
carrier ceases to instantiate it; a successor realization preserves the
registered intervention--response role; and perturbing that successor changes
the future family under matched controls.
\end{quote}
Biology is the natural place to test this statement, because component
turnover with persisting function is ubiquitous there and carrier
interventions are experimentally routine.

\section{The adjacent possible and its carriers}
\label{sec:ap}

Of the native formalisms, the adjacent possible is the closest to ICP and the
one in which ICP has the clearest opportunity to contribute. This section sets
out what the formal theory already establishes, identifies the step it leaves
unspecified, and states a prediction that separates a carrier-based account
from the existing models.

\subsection{What is formalized}

Kauffman defined the chemical adjacent possible as the molecular species absent
from the actual but synthesizable from it in one reaction step
\citep[p.~142]{Kauffman2000Investigations}, and later argued in this journal
that enabling constraints open unprestatable possibilities in biological
evolution \citep{Kauffman2014Prolegomenon,PatraKauffman2025NoLaw}. Tria,
Loreto, Servedio, and Strogatz gave the idea its first minimal mathematical
form \citep{TriaEtAl2014Novelties}. In their urn model with triggering, an
element drawn from an urn is returned together with $\rho$ copies
(reinforcement); if and only if the element is new, $\nu+1$ brand-new elements
are also added, so that $\nu+1$ is the size of the adjacent possible opened by
that novelty. The balance of the two processes fixes the statistics: the number
of distinct elements grows as $D(N)\sim N^{\nu/\rho}$ for $\nu<\rho$, as $N\log
N$ for $\nu=\rho$, and linearly for $\nu>\rho$, with Zipf's law $f(R)\sim
R^{-\rho/\nu}$. A semantic-triggering variant reproduces the clustering of
related novelties observed in Wikipedia edits, music listening, annotation tags,
and texts \citep{TriaEtAl2014Novelties}.

The framework has since been extended in several directions. Explorations of
the adjacent possible in social contact space reproduce the rate at which
individuals acquire new acquaintances and the degree, clustering, and community
structure of real networks \citep{UbaldiEtAl2021SocialNetworks}. Waves of
novelties in which newcomers are not suppressed by established elements have
been quantified across several systems \citep{MonechiEtAl2017Waves}. A
time-dependent urn with triggering, with reinforcement and triggering
parameters that vary in time, reproduces the logarithmic novelty growth
observed in gene mutations and patent codes and has a critical region in which
Heaps' and Zipf's laws coexist \citep{BellinaEtAl2025TimeDependentUrn}.
Edge-reinforced walks place discovery on an explicit concept network
\citep{IacopiniEtAl2018Innovation}, and higher-order novelty counts first
occurrences of combinations, distinguishing sequences with similar first-order
discovery rates \citep{DiBonaEtAl2025HigherOrder}. Combinatorial TAP models
generate long plateaus followed by explosive growth
\citep{SteelEtAl2020Combinatorial}, while long-run product data show
resource-constrained rather than unrestricted combinatorial discovery
\citep{Taalbi2026AdjacentPossible}. Fari\~na and Vallejo give a geometric
formulation, a biological cosmological constant proportional to the squared
fractional expansion rate of Fisher-information volume
\citep{FarinaVallejo2026LambdaB}.

ICP therefore cannot claim history-dependent accessibility, triggered support
growth, correlated novelty, a plateau-to-explosion transition, the Heaps and
Zipf regularities, or a geometric measure of expansion. All of these are prior
art.

\subsection{What is left unspecified}

In each of these models the opening step is postulated. The triggering
parameter $\nu$, or its time-dependent generalization, is fitted to data, and
the elements it adds are featureless new colours or nodes; what in the system
produces them lies outside the model. Two further features follow from this
construction. First, the basic urn never removes elements, so possibilities
once opened stay open indefinitely: in the ledger of
Eq.~\eqref{eq:possibility-flux-ledger} only the enabled set $E$ is populated,
and the lost set $L$ is empty by construction. Second, the opened
possibilities depend on the history only through the count and order of past
novelties, not through any structure that remains physically present.

Kauffman's enabling constraints name the missing element without making it
operational. This is the step ICP's carrier criterion is designed for. In its
terms, an account of how the adjacent possible expands must identify an
embodied carrier $m$, formed by the novelty event, whose presence makes the
frontier $\mathcal J^{\pi,\ell}_{t,\tau}$ of Eq.~\eqref{eq:adjacent-frontier}
larger than it is under a matched reference $m_0$. The triggering size then
becomes an intervention--response quantity,
\begin{equation}
 \nu_{\mathrm{C}}(m)
 :=\bigl|\mathcal J^{\pi,\ell}_{t,\tau}(m)\bigr|
  -\bigl|\mathcal J^{\pi,\ell}_{t,\tau}(m_0)\bigr|,
 \label{eq:carrier-triggering}
\end{equation}
measured by ablating, restoring, or replacing the carrier rather than fitted to
a sequence, and reported together with the identities of the enabled and lost
types. ICP does not supply the mechanism that makes $\nu_{\mathrm{C}}$ positive
in a given system; that is chemistry, biophysics, or sociology. It supplies the
test that would identify the mechanism, and the accounting under which the
identification can fail.

\subsection{A discriminating prediction}

The two accounts differ on one point that can be tested. In an urn model,
possibilities opened by a novelty persist whatever happens afterwards. On the
carrier account, an opened possibility persists only while its carrier, or a
role-equivalent successor, remains realized:
\begin{quote}
\textbf{Carrier dependence of the adjacent possible.} If the carrier formed by
a novelty event is removed after the event, without a successor, the
possibilities it opened leave the frontier, and they return if the carrier is
restored.
\end{quote}
A pure triggering model predicts no change in the frontier under this
intervention, because the history of novelties is unchanged. A complete
mechanistic model of a particular system may of course predict the same as the
carrier account; the prediction distinguishes ICP from the phenomenological
models of the adjacent possible, not from complete-state dynamics.

Molecular evolution already provides a case in which the prediction holds.
During the evolution of cortisol specificity in the ancestral glucocorticoid
receptor, permissive substitutions were functionally neutral in isolation but
had to be present for the large-effect function-switching mutations to be
tolerated \citep{HarmsThornton2014Contingency}. In the terms above, the
permissive substitutions are the carrier, the set of tolerable
function-switching mutations is the enabled set $E$, and the dependence of
tolerance on their presence is the carrier dependence of that frontier. The
same study found that permissive genotypes were extremely rare: a library of
thousands of variants of the ancestral protein recovered the historical
permissive substitutions but no alternatives, because a permissive mutation
must stabilize specific local structure, preserve the energetic balance between
conformations, and be compatible with both ancestral and derived structures
\citep{HarmsThornton2014Contingency}. Contingency in which few paths are open
also appears in the exhaustive reconstruction of TEM $\beta$-lactamase
trajectories \citep{WeinreichEtAl2006MutationalPaths}. These are cases in which
the opener of an adjacent possible has been identified physically and in which
its ablation is experimentally meaningful.

The pattern also connects to the crossover of
Eq.~\eqref{eq:constraint-enablement-crossover}. Permissive substitutions
restrict nothing measurable at the level of immediate receptor function, yet
they enlarge the set of later functional changes; a carrier that prunes
microstates while enlarging the role frontier is the general form of the same
observation.

\subsection{From statistics to carriers}

The urn models explain the statistics of discovery from a small number of
parameters. The carrier account asks what sets those parameters. If it is
correct, the triggering and reinforcement rates of a system should be
predictable from measured properties of its carriers: how often novelty events
form a carrier, how large a frontier each carrier opens, and how long carriers
survive or are succeeded. A time-dependent triggering rate
\citep{BellinaEtAl2025TimeDependentUrn} would then be read as a changing
population of realized carriers. Making that link quantitative in one system,
by predicting $\nu$ and its time dependence from carrier measurements before
fitting the sequence, is the most direct way in which ICP could become an
explanation of the adjacent possible rather than a redescription of it.

\section{Biological records and their ordinary rivals}
\label{sec:biology}

Biology already supplies the architecture
\begin{equation}
 \text{history}\rightarrow\text{embodied record}
 \rightarrow\text{changed later response},
 \label{eq:bio-architecture}
\end{equation}
including cases in which the initiating event or molecule is gone. In bacterial
acquired immunity, phage exposure installs a sequence record in a CRISPR locus,
and adding or deleting matching spacers changes later resistance
\citep{BarrangouEtAl2007CRISPR}. Exhaustive construction of all 32
combinations of five TEM $\beta$-lactamase substitutions found only 18 of 120
direct mutational paths accessible under the measured landscape
\citep{WeinreichEtAl2006MutationalPaths}, and ancestral reconstruction of the
glucocorticoid receptor showed that rare permissive substitutions made later
functional changes tolerable \citep{HarmsThornton2014Contingency}: historical
contingency acting on operational accessibility, a hard and an operational
constraint in the terms of Section~\ref{sec:formalism}. Constructive
origin-of-life work supplies controls in both directions. Cooperative RNA
networks can outgrow selfish cycles \citep{VaidyaEtAl2012CooperativeRNA}, while
compositional autocatalytic networks show that self-maintenance need not yield
selected heredity \citep{VasasEtAl2010Evolvability}.

Constraint-closure theory gives these cases a common organizational reading:
an enzyme may constrain one reaction while being produced and degraded on
another timescale, so constraints persist through the turnover of their
realizers \citep{MontevilMossio2015Closure}. Whether closure grounds function or
normativity remains contested \citep{Garson2017AgainstOrganizational,
Corti2023OrganizationalNormativity}.

Every example nonetheless has a sufficient ordinary mechanism: sequence and
base pairing, chromatin chemistry, protein structure, population dynamics, and
reaction kinetics. A state description that omits the spacer, mark, or enzyme
will appear history-dependent; an adequate physical state includes the
carrier. Generic future-shaping is therefore not ICP novelty. What biology does
provide is experimentally controllable systems in which the full criterion of
Section~\ref{sec:criterion}---including succession---can be applied and,
crucially, can fail. The next section applies it to one such system in detail.

\section{Worked case: transcriptional memory through chromatin succession}
\label{sec:ino1}

\subsection{The system}

Budding yeast cells previously starved of inositol reactivate the \textit{INO1}
gene faster on a second starvation than cells without that history. The memory
requires a cis-acting DNA element (the memory recruitment sequence, MRS), the
transcription factor Sfl1 bound to it, interaction with the nuclear pore
complex, incorporation of H2A.Z, and a memory-specific dimethylation of histone
H3 lysine 4 (H3K4me2) \citep{LightEtAl2010INO1,SumpEtAl2022INO1Memory}. The mark
is written by a form of COMPASS lacking the Spp1 subunit and requires SET3C, a
complex that binds H3K4me2 and interacts physically with the writer; the
authors propose this writer--reader coupling as the mechanism for
re-incorporating the mark after DNA replication \citep{SumpEtAl2022INO1Memory}.

Sump and colleagues separated establishment from inheritance by inducible
degradation of Sfl1. When Sfl1 was degraded after memory had been established,
peripheral localization was lost but H3K4me2 remained at the promoter for
roughly three to four generations; degrading Set3 together with Sfl1 caused
rapid loss of the mark \citep{SumpEtAl2022INO1Memory}. The authors conclude
that memory duration is governed by the number of divisions or replication
cycles, reflecting dilution, degradation, or imperfect fidelity of inheritance.

\subsection{Mapping onto the criterion}

Each replication fork distributes parental nucleosomes between the daughter
chromatids and incorporates new ones, so the physical realizations
$c_0,c_1,c_2,\ldots$ of the mark are progressively disjoint in material
composition. This is the situation the succession clause of
Section~\ref{sec:criterion} was written for. Table~\ref{tab:ino1} gives the
mapping.

\begin{table}[ht]
\centering
\small
\caption{The persistent-carrier criterion applied to \textit{INO1}
transcriptional memory. ``Shown'' refers to published results
\citep{LightEtAl2010INO1,SumpEtAl2022INO1Memory}; ``untested'' marks components
of the criterion that no published experiment yet addresses.}
\label{tab:ino1}
\begin{tabularx}{\linewidth}{>{\raggedright\arraybackslash}p{0.2\linewidth}
  >{\raggedright\arraybackslash}X
  >{\raggedright\arraybackslash}p{0.13\linewidth}}
\toprule
Element & Instantiation & Status \\
\midrule
History $m$ vs.\ $m_0$ & Prior inositol starvation vs.\ never-starved cells in the
same current medium & Shown \\
Future family $\Efamily$ & Distribution of \textit{INO1} re-induction times
after re-starvation, indexed by the delay $\pi\in\Pi$ & Shown (faster
reactivation) \\
Constraint type & Soft: naive cells also induce, so the support is unchanged
and the weights shift; Eq.~\eqref{eq:icp-constraint-divergence} applies
directly & Inferred \\
Formation & Sfl1 and pore interaction required to establish the mark & Shown \\
Carrier intervention & Removing writer/reader components (COMPASS, SET3C)
abolishes the mark and faster reactivation; matched shams required & Shown,
shams partial \\
Succession & H3K4me2 re-incorporated onto new nucleosomes for 3--4 generations
after initiator removal & Shown for the mark \\
Role equivalence $\sim$ & Successor cells whose re-induction distribution lies
within a declared tolerance of the memory distribution & Untested across
succession \\
Role exhaustion & Mark lost after finitely many divisions while the lineage
lives: $\Rset_C=\varnothing$ for the role, not the substrate & Shown for the
mark \\
\bottomrule
\end{tabularx}
\end{table}

\subsection{What the case shows and does not}

The case passes formation, sufficiency, and carrier intervention for the mark,
and it demonstrates succession of the \emph{record} through several rounds of
replication-coupled carrier replacement. It does not yet demonstrate
succession of the \emph{role}. Sfl1 is itself required to express the faster
reactivation, so while Sfl1 is absent the future-response family cannot be
assayed: the inherited mark is a latent record whose functional reading has not
been measured after succession. This distinction between persistence of a
record and persistence of a role is exactly the one ICP exists to enforce, and
the published literature does not need ICP to draw it. It does, however, make
the missing experiment precise.

\paragraph{Prospective test.} Establish memory, degrade Sfl1, allow
$k=0,1,\ldots,5$ generations of latent inheritance, restore Sfl1, and measure
re-induction after re-starvation. ICP's succession clause predicts that the
divergence $\mathfrak C_k$ between memory and naive re-induction-time
distributions decays with $k$ at the per-generation rate of H3K4me2 retention at
the locus, and that degrading Set3 only during the latent interval abolishes
$\mathfrak C_k$ while a matched sham degradation does not. If restored Sfl1
re-establishes fast reactivation independently of the inherited mark, the
role was not carried by the successor record and the case fails the
succession gate.

\paragraph{The ordinary rival.} Nucleosome-modification models with recruited
positive feedback already predict that the stability of a chromatin memory is
set by the balance between feedback-driven re-modification and
replication-coupled dilution, with bistable, effectively heritable states above
a feedback threshold \citep{DoddEtAl2007EpigeneticMemory}. A complete chromatin
model of this kind, with the writer and reader machinery included, fixes the
future family in the sense of Eq.~\eqref{eq:complete-state-null}. We expect
it to predict the outcome of the prospective test, and we do not claim that ICP
explains \textit{INO1} memory better than chromatin dynamics does.

The case is therefore a positive control for the criterion rather than
evidence for a new mechanism. A criterion that failed to classify a known,
mechanistically understood memory correctly would be disqualified; one that
classifies it correctly, separates record from role, and specifies the
untested experiment has earned the right to be applied where the mechanism is
not known. The same table can be filled in for CRISPR spacers, prion
conformational templates, and artificial systems, which makes them comparable
as instances of one question.

\section{An artificial-life testbed}
\label{sec:testbed}

The companion study applies the same criterion to a stochastic finite local
lattice model used as an artificial-life testbed \citep{Clark2026IotaConstraint}.
Its central result concerns a relational connectivity role rather than a
marked payload: in three prospectively replayed events across two seeds,
tracked regions lost connectivity and regained it through a causally necessary
successor bridge sharing no cells with its predecessor, and removing the
successor destroyed the restored relation. Event-resolved vetoes further showed
that alternative successors are frequently available but that an available
future is causally relevant without being sufficient. Within that model
Eq.~\eqref{eq:iota-kl-identity} holds exactly for its support measure. These
results establish possibility within one engineered law, not use by living
systems; they illustrate hard constraints, complementing the soft constraint of
Section~\ref{sec:ino1}.

\section{Adversarial audit and verdict}
\label{sec:audit}

\subsection{The complete-state challenge}

Let $Z_t$ be a complete physical specification at time $t$, including all
degrees of freedom relevant at the chosen precision, their correlations,
boundary conditions, resources, controls, and dynamics. Then
\begin{equation}
  \Efamily=\Phi_{\tau}(Z_t;\Pi,\mathcal A,B,C).
  \label{eq:complete-state-null}
\end{equation}
On this account a memory, chromatin mark, or regulatory state matters only
because it is part of $Z_t$; enlarging the state until it is sufficient removes
any independent history variable. This null is not defeated by showing that two
visibly similar systems behave differently, by finding mutual information
between past and future, or by a compact description alone. ICP becomes
distinct only if its carrier-equivalence construction does useful work beyond
Eq.~\eqref{eq:complete-state-null}: a minimal manipulable sufficient class, an
invariant across genuine succession, a validated compression of
intervention--response laws, or a preregistered prediction missed by the best
complete-state model. The \textit{INO1} case shows the criterion operating
correctly; it does not supply such a result.

\subsection{Failure modes that remain live}

\begin{enumerate}
  \item \textbf{State-augmentation collapse.} A claimed historical effect
  disappears when the carrier or environment is adequately represented,
  leaving ICP as vocabulary for ordinary memory.
  \item \textbf{Representation dependence.} The carrier exists only under a
  convenient partition, encoding, or coarse-graining.
  \item \textbf{Inferential--physical equivocation.} An observer's loss of
  uncertainty is mistaken for a change in the system's possibilities.
  \item \textbf{Generic-dynamics redescription.} Without the formation,
  intervention, succession, and robustness gates, ICP collapses into
  reachability, viability, or semantic-information theory.
  \item \textbf{Model artifact.} Effects in an artificial system depend on one
  update rule, scheduler, or topology.
  \item \textbf{Novelty by label.} A first visit to a prestated state or
  analyst-defined type is mistaken for creation of a new possibility.
  \item \textbf{Self-confirming evidence.} Project papers and pipelines share
  assumptions and error modes; external replication is required.
\end{enumerate}

\subsection{Graded verdict}

\begin{table}[ht]
\centering
\small
\caption{Graded verdict. Ratings refer to ICP as a research principle, not to
the truth of any ontology.}
\label{tab:verdict}
\begin{tabularx}{\linewidth}{>{\raggedright\arraybackslash}p{0.21\linewidth}
  >{\raggedright\arraybackslash}p{0.15\linewidth}
  >{\raggedright\arraybackslash}X}
\toprule
Dimension & Rating & Evidential basis and ceiling \\
\midrule
Conceptual coherence & Strong & Formation, intervention, future family, and
succession are separable and operationally specifiable. \\
Operational precision & Moderate--strong & The carrier gates and matched
contrasts can fail; the signed ledger handles pruning and enabling; no accepted
invariant yet exists. \\
Distinctness from prior art & Partial & Succession beyond the original carrier
is narrower than reachability or adjacent-possible expansion, but
complete-state, causal-state, and semantic-information accounts remain live. \\
Biological applicability & Moderate & The criterion classifies a known
chromatin memory correctly and identifies an untested role-succession
experiment. \\
Project evidence & Bounded & Artificial-life interventions give constructive
examples, not observations of living systems. \\
Falsifiability & Strong locally & Carrier edits, state augmentation,
representation changes, and held-out interventions can reject individual
claims. \\
\bottomrule
\end{tabularx}
\end{table}

\begin{quote}
\textbf{ICP is a coherent, falsifiable operational hypothesis about persistent
organization. Existing science supplies strong precedents for each of its
components but not for their integration. It is a viable research principle,
not an established law.}
\end{quote}

\section{Research path and conclusion}
\label{sec:path}

The program should advance in an order in which each step can fail.

\paragraph{Role succession in a living system.} Carry out the prospective
\textit{INO1} test of Section~\ref{sec:ino1}, and its analogues in other
systems with experimentally separable record and role: prion conformational
inheritance, CRISPR spacer acquisition and loss, and enzyme networks with
measured turnover. Freeze the role equivalence, horizon, and shams before
opening outcomes, and fit the best ordinary model on training interventions
before evaluating both on held-out ones.

\paragraph{Replication across artificial laws.} Repeat the testbed's
succession protocols under alternative local laws, schedulers, topologies, and
independent implementations, requiring
Eq.~\eqref{eq:icp-constraint-divergence} to keep its qualitative result.

\paragraph{Adjacent-frontier assay.} On each system, freeze the type map,
actualized repertoire, policies, and horizons; report enabled, lost, and
retained identities (Eq.~\eqref{eq:possibility-flux-ledger}) rather than novelty
counts; and compare against fixed-state, reinforced-urn, branching, and TAP
rivals. Test the carrier dependence of the frontier
(Section~\ref{sec:ap}) by removing and restoring a carrier after its novelty
event, and attempt to predict a system's triggering rate from carrier
measurements before fitting its sequence. The strongest near-term signature is a prospectively predicted
crossover of the form of Eq.~\eqref{eq:constraint-enablement-crossover}, with
$h_\star$ forecast from separately measured carrier lifetimes.

\paragraph{Stopping rules.} The principle should be narrowed or reformulated
if complete-state augmentation repeatedly removes the candidate historical
variable, if carrier interventions fail under matched controls, if the role
disappears under reasonable representation changes, or if no compression,
invariant, prohibition, or prediction emerges beyond native theories. A result
that reduces ICP to a precise theory of physically embodied memory would still
be valuable.

\paragraph{Conclusion.} Information can persist beyond the object that first
carried it, but every established example does so through continuing physical
realization in another state, relation, or organization. Biology shows records
outliving their initiators, constraints maintained through the turnover of
their realizers, and futures opened by prior organization. ICP's contribution
is a single operational test, with a signed formalism for pruning and enabling,
that makes these cases comparable and that separates the persistence of a
record from the persistence of a role. Its claim to be more than an organizing
vocabulary rests on a result it has not yet produced: a held-out prediction,
from carrier succession, that the best complete-state model of the same system
misses.



\section*{Data and code availability}
The literature registry, exact search log, claim--evidence matrix, rival and
falsification registers, formal derivations, and manuscript sources are
available at \url{https://github.com/fmertz55/ICP}. The simulation code,
preregistrations, and frozen outputs of the companion artificial-life study are
available at \url{https://github.com/fmertz55/iota}.

\section*{Ethics statement}
This theoretical and computational study used no human participants, human
tissue, animals, or field samples; ethics approval and consent were not
applicable.

\section*{Funding}
This research received no external funding.

\section*{Declaration of competing interests}
The author declares no known competing financial interests or personal
relationships that could have appeared to influence the work reported here.

\section*{CRediT author statement}
Douglas Clark: Conceptualization; Methodology; Formal analysis; Investigation;
Data curation; Software; Validation; Visualization; Writing -- original draft;
Writing -- review and editing; Project administration.

\section*{Declaration of generative AI and AI-assisted technologies}
During preparation of this work, the author used OpenAI Codex and Anthropic Claude for
AI-assisted literature triage, drafting, code review, and LaTeX editing. The
author reviewed and verified the cited sources, calculations, code, claims, and
final text and takes full responsibility for the content of the publication.

\bibliographystyle{plainnat}
\bibliography{library/references}

\end{document}
